\documentclass[11pt,a4paper]{article}

\usepackage[T1]{fontenc}
\usepackage[utf8]{inputenc}
\usepackage{lmodern}
\usepackage{amsmath,amssymb}
\usepackage{booktabs}
\usepackage{array}
\usepackage{hyperref}
\usepackage{xurl}
\usepackage{geometry}
\usepackage{enumitem}
\usepackage{listings}
\usepackage{caption}
\usepackage{microtype}

\geometry{margin=1in}
\hypersetup{
  colorlinks=true,
  linkcolor=blue,
  citecolor=blue,
  urlcolor=blue,
  pdftitle={School Van Tracking: Ordering Guarantees Under Unreliable Networks},
  pdfauthor={Anand Theertha}
}

\lstset{
  basicstyle=\ttfamily\small,
  breaklines=true,
  frame=single,
  columns=fullflexible,
  keepspaces=true
}

\title{School Van Tracking: Ordering Guarantees Under Unreliable Networks}
\author{Anand Theertha\\[0.5em]\normalsize Independent researcher}
\date{November 2026}

\begin{document}
\maketitle

\begin{abstract}
We study the telemetry protocol of a real-time school-van tracking application in which a driver's phone streams GPS positions to parents over unreliable mobile networks. Using TLA+ and the TLC model checker, we model \emph{two} unordered channels (driver$\rightarrow$server and server$\rightarrow$parent) with loss on the uplink and unbounded redelivery of the same sequence number---receives do not remove frames from the channel sets. We check that the position shown to a parent never moves backward (\texttt{MonotonicDisplay}). TLC reports a violation in the unguarded baseline; a seven-state counterexample is documented with raw logs in the public artifact. An ablation study shows that, under this downlink model, \emph{the parent display guard alone is sufficient} to restore monotonicity: a server-only guard still allows the map to jump backward when frames~1 and~2 both sit on the parent channel. Production deploys \emph{both} guards so server state and ancillary consumers (e.g., trip history) stay consistent, not because the server guard alone secures the live map. We also report fault-injection tests on the implementation (duplicate, reordered, and batched frames). We state precisely what the model does and does not prove and outline follow-on work (session epoch, shared acknowledgment state, retransmission).
\end{abstract}

\noindent\textbf{Index Terms---}formal verification, TLA+, model checking, GPS telemetry, mobile networks, message ordering, transportation systems.

\section{Introduction}
\label{sec:intro}

Parents who track a school van need a map position that reflects the vehicle's recent path. If the displayed position jumps \emph{backward} along the route---because an older GPS frame is applied after a newer one---they may infer the wrong direction, stop, or miss a pickup window. Cellular dead zones, reconnect bursts, and multi-node ingestion make such orderings plausible even when each TCP or WebSocket connection is in-order locally.

This paper is an applied case study: we (1)~formalize a minimal telemetry protocol and channel assumptions in TLA+, (2)~use TLC to find a monotonicity violation and to compare server-side versus client-side sequence guards, and (3)~relate the model to a deployed Node.js/React Native stack through fault injection. Our contributions are:
\begin{itemize}[noitemsep]
  \item A two-channel TLA+ model (lossy uplink, unordered downlink, redelivery) that abstracts reconnect and multi-publisher reordering as a channel assumption, with reproducible TLC logs.
  \item A guard \emph{ablation}: under the stated downlink, the parent display guard alone restores \texttt{MonotonicDisplay}; server-only guarding does not.
  \item A concise mapping from abstract violations to production mitigations and known remaining gaps (epoch on driver restart, cross-node races, liveness without acks).
\end{itemize}

Artifact: \url{https://github.com/Anand-Theertha/school-van-formal-verification} (release v1.2.0); archived copy DOI \url{https://doi.org/10.5281/zenodo.23567441}.

\section{System overview}
\label{sec:system}

The driver app samples GPS every two seconds and sends each position as a frame with trip id, monotonic \texttt{seq\_id}, timestamp, and coordinates:
\begin{enumerate}[noitemsep]
  \item Driver app (React Native) sends frames over Socket.IO on secure WebSockets.
  \item One or more Node.js ingestion nodes validate the token, persist to PostgreSQL/PostGIS, update a current-trip cache, and publish via a Redis pub/sub adapter.
  \item Parent apps subscribed to the trip receive the broadcast and render the position.
\end{enumerate}

The update rate is bounded by the sampling interval; measured average end-to-end latency was about 250\,ms on 4G in testing. Appendix~\ref{app:arch} summarizes stack, schema, and security choices.

\section{Baseline protocol and failure model}
\label{sec:failures}

\textbf{Baseline protocol.} The server accepts each valid frame and broadcasts it. Without sequence guards, the parent display follows whatever arrives on the downlink.

\textbf{Failure model.}
\begin{itemize}[noitemsep]
  \item Connections drop and reconnect; buffered frames may arrive in bursts.
  \item Application-level retries can duplicate deliveries.
  \item Frames from an old socket can arrive after frames from a new connection.
  \item Several server nodes can process frames for the same trip concurrently.
\end{itemize}

\textbf{Out of scope.} GPS error, dead reckoning, map matching, and malicious clients.

\section{Formal model (TLA+)}
\label{sec:model}

\subsection{Channel abstraction}

TCP and WebSockets do not reorder frames within one connection. We nevertheless use \emph{unordered} channels because that is a conservative abstraction of overlapping old and new sockets during reconnect, application retries, and parallel processing on multiple ingestion nodes.

We use two channels:
\begin{itemize}[noitemsep]
  \item \textbf{Uplink} (\texttt{network\_buffer}): driver$\rightarrow$server; \textbf{lossy} (\texttt{NetworkDrop}).
  \item \textbf{Downlink} (\texttt{parent\_channel}): server$\rightarrow$parent; \textbf{unordered} in the model.
\end{itemize}

\textbf{Deployment note.} Redis pub/sub may preserve order from a \emph{single} publisher. We model an unordered downlink because multiple ingestion nodes can publish toward the same trip, so parents need not see frames in send order. Multiple publishers are not modeled explicitly: there is a single server with one shared \texttt{server\_ack\_seq}, so downlink reordering is an assumption of the channel rather than a consequence of modeled nodes.

\textbf{Duplication.} Each channel is a set of sequence numbers; receives do \emph{not} remove an element, so the same \texttt{seq\_id} can be delivered arbitrarily many times (redelivery). We do not model two distinct in-flight copies of the same \texttt{seq\_id} in one set (a multiset would be needed for that).

A counterexample in this model shows the protocol \emph{abstract} can violate a property under these assumptions; it is not a proof about the live cellular network or the full production codebase.

\subsection{State variables}

\begin{center}
\begin{tabular}{@{}ll@{}}
\toprule
Variable & Meaning \\
\midrule
\texttt{driver\_seq} & Latest sequence number produced by the driver \\
\texttt{network\_buffer} & Set of sequence numbers in flight on the uplink \\
\texttt{server\_ack\_seq} & Highest sequence number the server has accepted \\
\texttt{parent\_channel} & Set of sequence numbers forwarded toward the parent \\
\texttt{parent\_view\_seq} & Sequence number currently displayed \\
\bottomrule
\end{tabular}
\end{center}

\subsection{Actions}

The specifications include exactly these actions (no separate ``connection drop'' or ``duplicate'' operators):

\begin{center}
\begin{tabular}{@{}l>{\raggedright\arraybackslash}p{0.68\textwidth}@{}}
\toprule
Action & Description \\
\midrule
\texttt{DriverSend} & Driver produces the next frame (\texttt{driver\_seq} increases, up to \texttt{MaxSeq}) \\
\texttt{NetworkDrop} & Uplink loss: remove one frame from \texttt{network\_buffer} \\
\texttt{ServerReceive} & Server reads one frame from \texttt{network\_buffer}; may append to \texttt{parent\_channel} \\
\texttt{ParentReceive} & Parent reads one frame from \texttt{parent\_channel}; updates \texttt{parent\_view\_seq} \\
\bottomrule
\end{tabular}
\end{center}

\textbf{Modules (ablation).}
\begin{itemize}[noitemsep]
  \item \texttt{SchoolVanTracker\_Baseline}: no guards.
  \item \texttt{SchoolVanTracker\_ServerGuarded}: server guard only.
  \item \texttt{SchoolVanTracker\_ParentGuarded}: parent guard only.
  \item \texttt{SchoolVanTracker\_Guarded}: both guards, as deployed.
\end{itemize}

\subsection{Bounds and specification}

\texttt{MaxSeq}~=~3 in all modules (finite state space for TLC). There is no bound on drops or redeliveries beyond finiteness of the model. \verb|Spec == Init /\ [][Next]_vars| with \textbf{no fairness} hypothesis for the safety runs below. TLC configurations set \texttt{CHECK\_DEADLOCK FALSE} (terminal states exist when the driver has finished and both channels are empty).

\section{Properties}
\label{sec:properties}

\textbf{Safety (action property):} the displayed position never moves backward.
\begin{lstlisting}
MonotonicDisplay == [][parent_view_seq' >= parent_view_seq]_parent_view_seq
\end{lstlisting}

\textbf{Safety (invariant):} only positions the driver has produced are displayed.
\begin{lstlisting}
NoFutureFrame == parent_view_seq <= driver_seq
\end{lstlisting}

\textbf{Liveness (not checked in current TLC configs):}
\begin{lstlisting}
FinalDelivery == (driver_seq = MaxSeq) ~> (parent_view_seq = MaxSeq)
\end{lstlisting}

\texttt{FinalDelivery} requires stronger fairness, bounded loss, or retransmission. Section~\ref{sec:limits} discusses acknowledgments for a future model revision. The TLC runs in this release check \texttt{MonotonicDisplay} and \texttt{NoFutureFrame} only.

\section{Model checking results}
\label{sec:results}

\textbf{Environment.} TLC~2.20 (\texttt{tla2tools.jar} via VS Code TLA+ extension), Temurin JDK~25, Windows~11, 2026-10-23. Raw logs ship in directory \texttt{tlc/} (Appendix~\ref{app:logs}).

\subsection{Baseline protocol}

TLC \textbf{violates} \path{MonotonicDisplay}. Shortest counterexample: \textbf{7 states} (six steps: two \texttt{DriverSend}, two server receives, two parent receives).

\begin{center}
\begin{tabular}{@{}cl@{}}
\toprule
State & Summary \\
\midrule
1 & Init \\
2--3 & Driver sends seq~1 and~2; \texttt{network\_buffer = \{1,2\}} \\
4--5 & Server accepts 1 then~2; \texttt{parent\_channel = \{1,2\}} \\
6 & Parent displays 2 (\texttt{parent\_view\_seq = 2}) \\
7 & Parent applies 1 $\rightarrow$ \texttt{parent\_view\_seq = 1} (\textbf{violation}) \\
\bottomrule
\end{tabular}
\end{center}

Reordering on the \emph{downlink} suffices; uplink duplication is not required for this trace. At error: 282 states generated, 122 distinct (search stopped with 64 states left on queue). Trace depth~7.

\subsection{Guard ablation (\texttt{MonotonicDisplay})}

\begin{center}
\begin{tabular}{@{}lllrrr@{}}
\toprule
Variant & Module & Result & Generated & Distinct & Depth \\
\midrule
No guards & Baseline & Violated & 282 & 122$^\dagger$ & 7 \\
Server only & ServerGuarded & Violated & 251 & 97$^\dagger$ & 7 \\
Parent only & ParentGuarded & Holds & 1726 & 347 & 11 \\
Both (v1) & Guarded & Holds & 922 & 199 & 11 \\
\bottomrule
\end{tabular}
\end{center}

\noindent$^\dagger$First violation stops exploration early; distinct counts are not the full state graph.

\textbf{Conclusion.} Under an unordered \texttt{parent\_channel}, the \textbf{parent guard alone is sufficient} for \texttt{MonotonicDisplay}; the server guard alone is not. The server-only failure follows directly from the channel assumption: the server forwards 1 then~2 in order, and the set-based \texttt{parent\_channel} redelivers~1 after the parent has shown~2. Both guards remain appropriate in production for server/history consistency and stale-frame filtering on the uplink.

``Holds'' means no violation within \texttt{MaxSeq = 3} and the stated channel model; it is not a proof for unbounded sequences or all real-network behaviors.

\subsection{Liveness}

\texttt{FinalDelivery} was not run in the bundled \texttt{.cfg} files. Informally, without driver-side retention and server acknowledgments, a lost final frame can leave the parent view stalled (aligned with injection test T6 in Section~\ref{sec:injection}). That is an explicit v1 trade-off: safety properties were the focus of model checking.

\section{Protocol revisions (production v1)}
\label{sec:fixes}

\textbf{Server filter.} Drop any frame whose sequence number is less than or equal to the server's current highest seen value for that trip (uplink guard).

\textbf{Parent guard.} Render a frame only if its sequence number is greater than the last one displayed.

TLC ablation (Section~\ref{sec:results}) shows the parent guard is sufficient for monotonic \emph{display} under the modeled downlink; the server guard supports consistent server state and consumers other than the live map.

\section{Implementation validation (fault injection)}
\label{sec:injection}

\begin{center}
\small
\begin{tabular}{@{}lllcc@{}}
\toprule
Test & Injected fault & Expected after fix & Baseline & v1 \\
\midrule
T1 & Old frame after newer & No backward jump & 14 jumps & 0 jumps \\
T2 & Duplicate frame & No visible change & UI flickers & Ignored \\
T3 & Burst after reconnect & Final position correct & Jumps & Correct \\
T4 & Driver restart, counter resets & Frames accepted & Accepted & Dropped (stale) \\
T5 & Two nodes concurrent & No backward jump & Jumps & Jumps (cross-node race) \\
T6 & Last frame lost before reconnect & Final position shown & Stalls & Stalls \\
\bottomrule
\end{tabular}
\end{center}

\section{Limitations and future work}
\label{sec:limits}

\begin{itemize}[noitemsep]
  \item \textbf{Session epoch} (T4): local \texttt{seq\_id} reset without epoch can freeze the trip.
  \item \textbf{Idempotent history:} store all frames with \((\text{trip}, \text{epoch}, \text{sequence})\); separate history from live broadcast.
  \item \textbf{Shared server state} (T5): atomic Redis update for highest seen seq across nodes.
  \item \textbf{Acknowledgments and retransmission} (T6): needed for \texttt{FinalDelivery}-style liveness.
  \item \textbf{Model:} optional multiset for duplicate in-flight copies; explicit downlink loss; per-action fairness for liveness checks.
\end{itemize}

\section{Related work}
\label{sec:related}

This work sits at the intersection of applied formal methods and distributed telemetry. Lamport's TLA+ and TLC toolchain~\cite{lamport2002} is the verification backbone. Kleppmann~\cite{kleppmann2017} discusses ordering, idempotency, and stream processing patterns relevant to guarded ingestion. Lamport's logical clocks~\cite{lamport1978} motivate sequence-based reasoning without claiming equivalence to our guards. Beschastnikh et al.~\cite{beschastnikh2016} survey debugging distributed systems; our counterexample traces play a similar explanatory role at design time.

\section{Conclusion}
\label{sec:conclusion}

Unreliable networks and multi-node publish paths can expose parents to non-monotonic map updates even when each connection is locally ordered. A small TLA+ model with explicit uplink loss and an unordered downlink yields a reproducible seven-state violation and a guard ablation that clarifies \emph{where} monotonicity must be enforced: at the parent display under our assumptions. Production adds a server guard for broader consistency; closing epoch, cross-node, and liveness gaps remains future work. Artifacts and TLC logs are public for independent replay.

\section*{Data and code availability}

TLA+ modules, TLC configuration files, unedited console logs (\path{tlc/}), modeling notes (\path{MODELING.md}), and this preprint source are available at \url{https://github.com/Anand-Theertha/school-van-formal-verification}. A versioned archive is deposited at Zenodo (concept DOI, resolving to the latest version: \url{https://doi.org/10.5281/zenodo.23567441}).

\begin{thebibliography}{9}

\bibitem{lamport2002}
L.~Lamport,
\emph{Specifying Systems: The TLA+ Language and Tools for Hardware and Software Engineers},
Addison-Wesley, 2002.

\bibitem{kleppmann2017}
M.~Kleppmann,
\emph{Designing Data-Intensive Applications},
O'Reilly, 2017.

\bibitem{lamport1978}
L.~Lamport,
``Time, clocks, and the ordering of events in a distributed system,''
\emph{Communications of the ACM}, vol.~21, no.~7, pp.~558--565, 1978.

\bibitem{beschastnikh2016}
I.~Beschastnikh, P.~Wang, Y.~Brun, and M.~D.~Ernst,
``Debugging distributed systems: Challenges and options for validation and debugging,''
\emph{Communications of the ACM}, vol.~59, no.~8, pp.~32--37, 2016.

\end{thebibliography}

\appendix

\section{Architecture and schema}
\label{app:arch}

\textbf{Stack.} Node.js/Express; Socket.IO over WSS with Redis pub/sub; PostgreSQL~15 with PostGIS; React Native driver and parent apps.

\textbf{Schema (conceptual).} Users, routes, children, trips, and location samples keyed by trip with monotonic \texttt{seq\_id}, geospatial columns (PostGIS), and time partitioning for retention.

\textbf{Security.} JWT (RS256); bcrypt password hashing; role-based publish/subscribe; TLS~1.3 on public endpoints.

\textbf{Design decisions.} WebSockets over polling for low latency; PostgreSQL/PostGIS over document stores for geospatial queries; shared ordering responsibility across driver buffering, server filtering, and parent display guards.

\section{TLC log excerpt (baseline violation)}
\label{app:logs}

The following lines are excerpted from \texttt{tlc/baseline-monotonic-2026-10-23.txt} (TLC~2.20, 2026-10-23). Full logs for baseline, server-guarded, parent-guarded, and both-guarded runs are in the repository.

\begin{lstlisting}[caption={Excerpt: error state (state 7).}]
7: <Next line 55, col 8 to line 55, col 49 of module SchoolVanTracker_Baseline>
/\ network_buffer = {1, 2}
/\ driver_seq = 2
/\ server_ack_seq = 2
/\ parent_view_seq = 1
/\ parent_channel = {1, 2}

282 states generated, 122 distinct states found, 64 states left on queue.
The depth of the complete state graph search is 7.
Finished in 8140ms at (2026-10-23 15:02:04)
\end{lstlisting}

\end{document}
